% Discussion. Interpretation, limitations and outlook; every number is \nm{key} from generated/numbers.tex.
\section{Discussion}
\label{sec:discussion}

A count matrix has been treated as a generic array: stored as coordinates or compressed columns, then handed to a
byte-level compressor that cannot see what the bytes mean. \sone{} instead codes what is known about counts. Most
non-zeros are ones; a cell's depth sets the rate of every gene in it; and genes ranked by how many cells express them
turn a cell's non-zeros into short gaps. With these regularities as contexts for a static entropy coder, the files are
\nm{ratio-h5ad_gzip}$\times$ smaller than gzip-compressed AnnData and \nm{ratioBestOther}$\times$ smaller than the most
compact general-purpose alternative, and they decode faster than any other format we measured, including uncompressed
arrays. The payload ladder (Fig.~\ref{fig:floor}a) shows where the gain comes from: a strong general-purpose compressor
applied to well-separated gap and value streams stops at \nm{ladder-L3_gaps_leb_zstd19} bits per non-zero, and the last
step to \nm{ladder-L4_s1_payload} comes from modelling, not from a better back-end coder. Genomic read formats learned
the same lesson~\cite{bonfield2022cram,numanagic2016comparison}.

\paragraph{Where the bytes are read.}
Size matters most where data move over a network. Atlases serve count matrices to thousands of users, and models that
train on them read every cell many times. A format that is both small and fast to decode moves the bottleneck back to the
model. Size alone is not the whole story, though: on a link with realistic latency, the number of round trips matters as
much as the number of bytes. The index at the head of every \sone{} member lets a reader plan all of a minibatch's
requests before sending any, and sending them together cut the time to a first minibatch by about
\nm{concurrency-gain-first}$\times$ on a 30\,ms link. The same plan-then-fetch pattern applies to any reader of
independently addressable chunks, and the \texttt{singlet} readers for atlas samples and for R use it as well.

\paragraph{How much is left.}
Cross-validated count models, including negative binomial and cluster-conditioned models, code the matrices a median of
\nm{s1AboveFloorPct}\% more compactly than \sone, and they are charged nothing for describing themselves. On read-count
protocols \sone{} codes the matrix more compactly than every model we fitted (Table~\ref{tab:floor}). The sampling bound
explains why the margin is small: about \nm{samplingPctMedian}\% of the \sone{} code length is the multinomial sampling
of molecules, which no lossless decoder can avoid. Knowing a cell's cluster is worth \nm{sharedBpnMedian} bits per
non-zero before the cost of the model. Larger gains therefore require giving up losslessness, for example by
quantising counts, or storing a model of the data rather than the data. Neither is acceptable for a primary archive,
and the second is outside the scope of a storage format.

\paragraph{Limitations.}
\begin{itemize}\itemsep 1pt
\item \emph{Hardware.} All timings come from \nm{threadsHeadline}--8 cores of shared cluster nodes of one type, with one of
  our timed jobs per node and other users' jobs beside it. Absolute rates depend on the CPU and disk; the comparisons
  between formats, measured on the same node and file, are what the conclusions rest on. On this CPU the AVX-512
  kernel gives no gain over AVX2.
\item \emph{Large matrices.} Above $4\times10^8$ non-zeros we timed a subset of formats, at one thread count and from
  local disk only. The R formats, RDS and BPCells, are not applicable above $2^{31}-1$ non-zeros, because an R
  \texttt{dgCMatrix} cannot hold such a matrix in memory. One planned dataset, a genome-wide Perturb-seq screen with
  6.3 billion non-zeros, was excluded because writing the comparison formats needed more memory and time than the
  campaign allowed.
\item \emph{Network tiers.} The network tiers are a local object store with a limit on total bandwidth and added
  latency; real object stores add variable latency and request costs. The 10\,Gbit/s tier exceeds what the test setup
  delivers and is flagged in the tables.
\item \emph{Scope.} \sone{} stores non-negative integer counts below $2^{32}$. Normalised or otherwise non-integer
  matrices are rejected, not rounded. Files are immutable; adding cells means writing a new sample object.
\item \emph{Writing.} We did not time writing. The \sone{} writer counts its tables in one pass over the matrix and then
  codes it, but the formats differ in how much writing costs.
\end{itemize}

\paragraph{Beyond single cells.}
Two out-of-domain count matrices, bulk RNA-seq gene counts and a bag-of-words corpus, were run as an extension. There
\sone{} is \nm{ext-gtex_bulk-ratio-h5ad_gzip}$\times$ and \nm{ext-nytimes_bow-ratio-h5ad_gzip}$\times$ smaller than
gzip-compressed AnnData. The same structure, heavy-tailed rank-frequency laws and many small counts, is what the codec
exploits, so any sparse count matrix of this kind can be stored the same way.
